\documentclass[11pt]{article}
\usepackage[T1]{fontenc}
\usepackage{amsmath,amssymb,geometry,hyperref}
\geometry{margin=1in}
\begin{document}
\begin{center}
{\Large Quasi-complete graphs need not maximize spectral moments}\\[4pt]
Disproof of Conjecture 00000003725\\[4pt]
ziangni-sys --- October 8, 2026
\end{center}
\paragraph{Moment convention and scope.}
For a finite simple graph $G$ with actual adjacency matrix $A_G$,
the standard adjacency spectral moment of order $k$ is
\[
 M_k(G)=\sum_j\lambda_j(A_G)^k=\operatorname{tr}(A_G^k).
\]
The trace equality follows by diagonalizing the real symmetric
adjacency matrix. It also counts length-$k$ closed walks.
We use $k=4$ and compare actual graphs with both vertex count and
edge count fixed. A single competing graph with larger fourth
moment is enough to refute the asserted quasi-complete extremum.
No lexicographic ordering of the entire moment sequence is stated
in the problem, and we make no claim about that different ordering.

\paragraph{The prescribed quasi-complete graph.}
The standard quasi-complete construction writes
$m=\binom{r}{2}+s$, $0\leq s<r$: take a clique on $r$ vertices,
join one new vertex to $s$ clique vertices, and leave any remaining
vertices isolated. This convention is described in
\emph{Sum of squares of degrees in a graph}, Section 1,
by Abrego, Fernandez-Merchant, Neubauer and Watkins (2009),
\url{https://emis.de/ft/14223}.
For $n=5,m=4$, the parameters are $r=3,s=1$.
Thus $Q=QC(5,4)$ has edges $01,02,12,03$ and isolated vertex $4$.
Compare it with the star $S=K_{1,4}$, whose edges are
$01,02,03,04$. Both are actual simple five-vertex, four-edge graphs.

\paragraph{Exact fourth moments.}
Their adjacency matrices, in vertex order $0,1,2,3,4$, are
\[
 A_Q=\begin{pmatrix}
 0&1&1&1&0\\1&0&1&0&0\\1&1&0&0&0\\1&0&0&0&0\\0&0&0&0&0
 \end{pmatrix},\qquad
 A_S=\begin{pmatrix}
 0&1&1&1&1\\1&0&0&0&0\\1&0&0&0&0\\1&0&0&0&0\\1&0&0&0&0
 \end{pmatrix}.
\]
Because the matrices are symmetric,
$\operatorname{tr}(A^4)=\sum_{i,j}(A^2_{ij})^2$.
For $Q$, the nonzero part of $A_Q^2$ is
\[
 \begin{pmatrix}3&1&1&0\\1&2&1&1\\1&1&2&1\\0&1&1&1\end{pmatrix},
\]
whose squared entries sum to $11+7+7+3=28$. For the star, $A_S^2$ consists of the central diagonal
entry $4$ and an all-one $4\times4$ leaf block, giving
$4^2+16=32$. Consequently
\[
 M_4(S)=32>28=M_4(Q).
\]
Hence $QC(5,4)$ does not maximize the actual fourth spectral
moment among graphs with the prescribed counts, and the conjecture
is false. The same strict comparison holds for moments divided by
the common vertex count five.

\paragraph{Formal certificate.}
Lean constructs the actual graphs and checks the canonical parameters
and both edge sets. It verifies the real adjacency matrices and
calculates the fourth-power traces. The final theorem disproves the
asserted maximum. All calculations are checked by the Lean kernel;
final audits use only standard Lean and Mathlib axioms.
\end{document}
